\documentclass[10pt,a4paper]{amsart}
\usepackage[margin=19mm]{geometry}
\usepackage{amsmath,amssymb,mathtools}
\usepackage[scr=rsfs,cal=euler]{mathalfa}
\usepackage{xfrac}
\usepackage[unicode,hidelinks,bookmarks=false]{hyperref}
\newcommand{\ie}{i.e.\ }
\newcommand{\cf}{cf.\ }
\newcommand{\resp}{respectively\ }
\newcommand{\Subsection}{\subsection}
\providecommand{\nobreakdash}{\nobreak\hbox{-}}
\setlength{\emergencystretch}{2em}
\multlinegap0pt
\usepackage[shortlabels]{enumitem}
\usepackage{mathtools}
\usepackage{siunitx}
\usepackage{float}
\usepackage{xcolor}
\usepackage{amssymb}
\usepackage{algorithm}\usepackage{algpseudocode}
\newtheorem{lemma}{Lemma}
\title{Scenario Reduction for Distributionally Robust Optimization}
\author{Kevin-Martin Aigner, Sebastian Denzler, Frauke Liers, Sebastian Pokutta, Kartikey Sharma}
\date{Source: arXiv:2503.11484v3 (CC BY 4.0)}
\begin{document}
\maketitle

\noindent\emph{Source and licence.} Scenario Reduction for Distributionally Robust Optimization. Version arXiv:2503.11484v3 (CC BY 4.0), distributed by arXiv under the Creative Commons Attribution 4.0 International licence, http://creativecommons.org/licenses/by/4.0/, as recorded in the arXiv licence metadata for this exact version and shown on the abstract page https://arxiv.org/abs/2503.11484v3. Full licence notice: \texttt{licenses/arxiv-2503.11484-CC-BY-4.0.txt}.\par\smallskip

\noindent\textit{Editorial scope.} Counted unit: the article's lemma lemma:ellipsoid\_transformation (source0.tex lines 2203--2210). The article's complete original proof of it is delivered with it (source0.tex lines 2211--2340, one \texttt{proof} environment of 130 lines). No mathematics is written, repaired or re-derived: the statement and the proof are the article's own lines, in the article's own notation, and no retained line is substituted or re-flowed.  Retained uncounted as the source-local context the proof needs: lines 2342--2345.  Nothing was omitted from any retained span, so the package carries no omission record.  No retained line contains a citation, so the wrapper carries no bibliography and no bibliography entry.  No retained line contains a figure environment, an image-inclusion command or drawing code, so no image is delivered and the package declares no asset dependency.  Editorial formatting only, in addition: an AMS \texttt{amsart} preamble that restates the article's own declarations and macros used by the retained lines, namely the environments \texttt{lemma} on separate counters, together with the standard packages those lines need.  The retained environments print in this document as Lemma 1 in order of appearance, whereas the article prints the same items with the numbering of its own full document.  The complete native source of the article as distributed in the arXiv e-print for this exact version is bundled unchanged as \texttt{sources/174767/source0.tex}.
\begin{lemma}
	\label{lemma:norm_bound}
	For any full row rank matrix \(A\), we have that \(\|A^{\top}(A A^{\top})^{-1}A\| \leq 1\) where \(\|A\| = \sup_{\|x\|_2 = 1} \|A x\|_2\). 
\end{lemma}

\begin{lemma}
	\label{lemma:ellipsoid_transformation}
	\begin{enumerate}[label=(\roman*)]
	\item The transformation of a sphere centered at the origin by a full row rank matrix \(A\) is an ellipsoid centered at the origin with covariance matrix \(A A^{\top}\). 
	
	\item The map of an ellipsoid $P \in \mathbb{R}^n$ centered at \(x_0\) and with convariance matrix \(\Sigma\) under a full row-rank linear transformation $A \in \mathbb{R}^{m \times n}$ with $n \geq m$ is an ellipsoid in $R^m$ centered at \(A x_0\) and with covariance matrix \(A \Sigma A^{\top}\).
	\end{enumerate}
\end{lemma}

\begin{proof}
	\emph{Part (i):}
	Consider a sphere \(S\) given by 
	\[
	S = \{x \mid x^{\top}x \leq \epsilon^2\}	
	\]
	Given a full row rank matrix \(A\), we want to show that the transformation of \(S\) by the matrix \(A\) is an ellipsoid given by 
	\[
	E = \{y \mid y^{\top}(A A^{\top})^{-1}y \leq \epsilon^2\}.	
	\]
	To do so, we need to show that \(A S \subseteq E\) and \(E \subseteq A S\). 

	\emph{First Case}.
	Consider a point \(x \in S\). We want to show that \(A x \in E\).
	To do so, we need to show that. 
	\[
	x^{\top}A^{\top}(A A^{\top})^{-1}A x \leq \epsilon^2.	
	\]
	We know that due to the Cauchy-Schwarz inequality 
	\[
		x^{\top}A^{\top}(A A^{\top})^{-1}A x \leq \|A^{\top}(A A^{\top})^{-1}A\| \|x\|_2^2	
	\]
	where the first norm is the matrix norm.
	%\todo{There are multiple different matrix norms, so we would have to write 'a' matrix norm.}
	Due to Lemma~\ref{lemma:norm_bound}, we observe that \(\|A^{\top}(A A^{\top})^{-1}A\| \leq 1\). 
	Thus
	\[
		x^{\top}A^{\top}(A A^{\top})^{-1}A x \leq \|x\|_2^2	\leq \epsilon^2,
	\]
	which shows that \(A S \subseteq E\). 
	
	\emph{Second Case}.
	Now, we want to show that \(E \subseteq A S\). For this, we consider \(y \in E\) and show that there exists a corresponding \(x \in S\) such that \(y = Ax\).
	We define \(x\) as \(A^{+}y\). Since \(A\) has linearly independent rows then \(A^{+} = A^{\top}(A A^{\top})^{-1}\).
	We can then show that 
	\[
	A x = A A^{+} y = y, 	
	\]
	because \(A^{+}\) will be a right inverse of \(A\). 

	We now show that \(x \in S\). 
	To do so, we need to show that \(x^{\top}x \leq \epsilon^2\).
	We can then write 
	\begin{align*}
		x^{\top}x &= (A^{+}y)^{\top}A^{+} y \\
		&= y^{\top}(A^{\top}(A A^{\top})^{-1})^{\top}A^{\top}(A A^{\top})^{-1} y\\
		&= y^{\top}((A A^{\top})^{-1})^{\top}A A^{\top}(A A^{\top})^{-1} y\\
		&= y^{\top}((A A^{\top})^{-1})^{\top} y\\
		&= y^{\top}(A A^{\top})^{-1} y\\
		&\leq \epsilon^2.
	\end{align*}
	Here, the second equality is due to the definition of the pseudo-inverse, the third due to the transpose, and the fourth because \(AA^{\top}(AA^{\top})^{-1} = I\). 
	The last equality is because \(A A^{\top}\) and its inverse would be symmetric. 
	Finally, the last inequality exists because \(y \in E\). 
	This concludes the proof. 

	\emph{Part (ii):}
	Let $P=\{x \in \mathbb{R}^n \; \vert \, (x-x_0)^T \Sigma^{-1} (x-x_0) \leq \epsilon^2 \}$ be an ellipsoid in $\mathbb{R}^n$ centered at $x_0 \in \mathbb{R}^n$ with a positive definite matrix $\Sigma \in \mathbb{R}^{n \times n}$.
	Let $A \in \mathbb{R}^{m \times n}$ be a full row rank matrix.

	To prove the result, we show that the ellipsoid \(P\) can be expressed as the translation and a linear transformation of a sphere.
	
	We perform a variable transformation $x' \coloneqq x-x_0$:
	\begin{equation*}
		P' = \{x' \in \mathbb{R}^n \; \vert \; x'^T \Sigma_x^{-1}x' \leq \epsilon^2\}
	\end{equation*}
	Then we can observe that \(P = \{x \;\vert \; x_0 + x',\; x' \in P'\}\).
	Here, \(P'\) is an ellipsoid centered at the origin. 

	Since $\Sigma$ is positive definite, there exists an orthogonal decomposition
	\begin{equation*}
		\Sigma = S D S^T
	\end{equation*}
	with $S^T S = 1_n$ and $D \in \mathbb{R}^{n \times n}$ diagonal.
	With the matrix $D^{-1/2}$ containing the inverse square roots of the diagonal entries, we calculate:
	\begin{equation*}
		x'^T \Sigma^{-1} x' = x'^T S D^{-1} S^T x' = x'^T S D^{-1/2} D^{-1/2} S^T x' = (D^{-1/2} S^T x')^T (D^{-1/2} S^T x').
	\end{equation*}
	We now define
	\begin{equation*}
		B \coloneqq D^{-1/2} S^T \text{ and} \quad x'' \coloneqq Bx',
	\end{equation*}
	which allows us to write
	\begin{equation*}
		x'^T \Sigma_x x' = (Bx')^T(Bx') = x''^T x''
	\end{equation*}
	This  we can define the set \(P''\) as 
	\begin{equation*}
		P'' \; \coloneqq \; \{x'' \in \mathbb{R}^n \vert x'' = Bx', \; x \in P' \}.
	\end{equation*}
	Due to the positive definiteness of $\Sigma$, we can invert $B$:
	\begin{equation*}
		x' = B^{-1}x'' = S D^{1/2} x''.
	\end{equation*}
	This shows that 
	\[
	P' = B^{-1} P''.	
	\]

	This, combined with our earlier result, allows us to write 
	\[
	P = \{x \;\vert\; x = x_0 + B^{-1}x'',\; x'' \in P''\}.
	\] 
	More generally, we can state
	\[P = x_0 + P' = x_0 + B^{-1}P''\].

	Now, we consider the linear transformation of \(P\).
	We then have 
	\[
	A P = \{y \;\vert\; y = A x_0 + A B^{-1}x'',\; x'' \in P''\}.	
	\]

	We define \(A' = A B^{-1}\). Since \(A\) is full row rank and \(B^{-1}\) is invertible, \(A B^{-1}\) has full row rank.
	Then, we can write
	\[
	A P = \{y \;\vert\; y = A x_0 + A' x'',\; x'' \in P''\}.	
	\]

	Due to Lemma~\ref{lemma:ellipsoid_transformation} and since \(P''\) is a sphere, the transformation of the sphere \(P''\) by a full row rank matrix \(A'\) is an ellipsoid with covariance matrix \(A' (A')^{\top}\).
	We have 
	\begin{align*}
		A' (A')^{\top} &= A B^{-1} (B^{-1})^{\top}A^{\top}\\
		&= A S D^{1/2} (S D^{1/2})^{\top}A^{\top}\\
		&= A S D^{1/2}  (D^{1/2})^{\top}S^{\top} A^{\top}\\
		&= A S D S^{\top} A^{\top}\\
		&= A \Sigma A^{\top}
	\end{align*}.

	Thus the linear transformation of \(P\) by a full row rank matrix \(A\) is an ellipsoid with center \(A x_0\) and covariance matrix \(A \Sigma A^{\top}\).
\end{proof}

\end{document}
